\documentclass[11pt]{article}
\usepackage[margin=1in]{geometry}
\usepackage{amsmath,amssymb,amsthm,mathtools,booktabs,microtype}
\usepackage[hidelinks]{hyperref}
\newtheorem{theorem}{Theorem}[section]
\newtheorem{proposition}[theorem]{Proposition}
\newtheorem{corollary}[theorem]{Corollary}
\newtheorem{remark}[theorem]{Remark}
\newcommand{\E}{\mathbb E}
\newcommand{\Pp}{\mathbb P}
\title{Dependent Gaussian Plaquette Polymers on a Four-Dimensional Gauge Block\\
\large Shared-link and Bianchi dependence without an independence hypothesis}
\author{}
\date{20 July 2026}
\begin{document}
\maketitle

\begin{abstract}
The independent-defect benchmark used in an ultraviolet Yang--Mills
calibration does not represent shared links or Bianchi constraints.  We replace
it by the exact linearized gauge quotient on the open side-two four-dimensional
block.  Its curvature is Gaussian with covariance $tP\otimes I_3$, where $P$
is the orthogonal projection onto the $136$-dimensional scalar exact-curvature
subspace.  Projection contraction yields, for every set of $k$ plaquettes, a
simultaneous radial-tail estimate $\exp[-kJ_3(\rho^2/t)]$, with
$J_3(s)=\tfrac12(s-3-3\log(s/3))$.  Combining this joint estimate with the
exact degree-$18$ adjacency graph gives a subcritical rooted connected-set sum
at the frozen ultraviolet witness, without assuming independence.  The result
is exact for the linearized Gaussian model; comparison with the nonlinear
$\mathrm{SU}(2)$ measure remains open.
\end{abstract}

\tableofcontents

\section{Gauge quotient in curvature coordinates}

Let $d_1:C^1\to C^2$ be the scalar link-to-plaquette coboundary.  On the
contractible block, $\ker d_1$ consists of gauge directions.  Quotienting that
kernel and mapping to curvature identifies the physical scalar space with
$\operatorname{im}d_1$, of dimension $136$.

For three color components, the Gaussian action is
\[
 S(F)=\frac1{2t}\|F\|^2,
 \qquad F\in\operatorname{im}d_1\otimes\mathbb R^3.
\]
The Jacobian of the linear quotient-to-curvature map is constant.  Therefore
$F$ is the isotropic Gaussian on this subspace.  In the ambient
$\mathbb R^{216}\otimes\mathbb R^3$ coordinates,
\begin{equation}
 \operatorname{Cov}(F)=tP\otimes I_3,
 \label{eq:cov}
\end{equation}
where $P$ is the orthogonal projection onto $\operatorname{im}d_1$.

The executable matrix has idempotence error $6.33\times10^{-15}$, eigenvalues
in numerical interval $[-7.04\times10^{-16},1+6.7\times10^{-15}]$, and rank
$136$.  These computations certify the declared finite matrix; the abstract
property $0\leq P\leq I$ follows exactly from orthogonal projection.

\section{A joint radial tail with dependence retained}

For a plaquette set $S$ of size $k$, let $F_S\in\mathbb R^{3k}$ be the
coordinate restriction.  Its covariance is a principal compression of
\eqref{eq:cov} and hence is bounded above by $tI_{3k}$.

\begin{theorem}[Dependent simultaneous-defect estimate]
Let $s=\rho^2/t$.  If $s>3$, then for every $S$ with $|S|=k$,
\begin{equation}
 \Pp\{\|F_p\|\geq\rho\text{ for all }p\in S\}
 \leq \exp[-kJ_3(s)],
 \qquad
 J_3(s)=\frac12\left(s-3-3\log\frac{s}{3}\right).
 \label{eq:joint}
\end{equation}
No independence assumption is made.
\end{theorem}
\begin{proof}
The event implies $\|F_S\|^2\geq k\rho^2$.  For
$0<\lambda<(2t)^{-1}$, covariance contraction and the Gaussian determinant
formula give
\[
 \E e^{\lambda\|F_S\|^2}
 \leq(1-2t\lambda)^{-3k/2}.
\]
Markov's inequality yields
\[
 \Pp\{\|F_S\|^2\geq k\rho^2\}
 \leq e^{-\lambda k\rho^2}(1-2t\lambda)^{-3k/2}.
\]
The optimizer is $2t\lambda=1-3/s$, which is admissible for $s>3$.
Substitution proves \eqref{eq:joint}.
\end{proof}

\section{Rooted connected-set theorem}

Use plaquette adjacency by a shared link.  Its exact maximum degree is
$\Delta=18$.  A canonical spanning-tree traversal bounds the number of
connected sets of size $k$ containing a fixed root by $\Delta^{2(k-1)}$.

\begin{corollary}[Dependent Gaussian connected-event sum]
Set $q_G=e^{-J_3(s)}$.  If
\begin{equation}
 r=e^\alpha\Delta^2q_G<1,
 \label{eq:ratio}
\end{equation}
then
\begin{equation}
 \sup_p\sum_{\substack{X\ni p\\ X\ \mathrm{connected}}}
 e^{\alpha|X|}\Pp\{X\text{ is entirely bad}\}
 \leq\frac{e^\alpha q_G}{1-r}.
 \label{eq:rooted}
\end{equation}
\end{corollary}
\begin{proof}
Apply \eqref{eq:joint} for each connected $X$, insert the connected-set count,
and sum the geometric series from $|X|=1$.
\end{proof}

\section{Ultraviolet witness}

For $t=3\times10^{-6}$ and $\rho=0.0082$,
\begin{center}
\begin{tabular}{@{}lr@{}}
\toprule
quantity & value\\
\midrule
$s=\rho^2/t$ & $22.4133333$\\
$J_3(s)$ & $6.69010107$\\
$q_G$ & $0.00124315712$\\
$e^{0.1}18^2q_G$ & $0.445143955$\\
rooted sum upper & $0.00247613973$\\
maximum certified $t$ & $3.27496193\times10^{-6}$\\
\bottomrule
\end{tabular}
\end{center}
The maximum heat time is obtained by monotone bisection of
$J_3(\rho^2/t)=\alpha+2\log18$.  The witness is strict, with heat-time slack
$2.74962\times10^{-7}$.

\section{Frontier audit}

This theorem shows that shared-link and linear Bianchi dependence do not by
themselves close the ultraviolet window.  It does not control nonlinear BCH
remainders, Haar densities, incoherent collars, noncentral topological data,
or repeated RG.  Equation \eqref{eq:rooted} is a connected bad-event envelope;
it is not by itself a Mayer or Ursell activity expansion.  The next bridge must compare the nonlinear measure with this
Gaussian law on a coherent collar and separately charge collar failure.  The
result is C2 for the linearized Gaussian gauge model.

\end{document}
